\chapter{Probability, conditioning, and estimation}
\label{ch:probability}

This chapter fixes the small amount of probability the rest of the document uses. A reader
comfortable with conditional expectation and the orthogonality principle can skip to
Chapter~\ref{ch:diffusion}, pausing only at \S\ref{sec:bayesopt}, whose result is used
constantly.

\section{Densities, joints, conditionals}

We work with random vectors on $\Reals^{\nsites}$ with densities. Write $P(a)$ for the density
of $a$, $P(a,x)$ for a joint density, and
\begin{equation}
P(a\mid x) \;=\; \frac{P(a,x)}{P(x)},
\qquad
P(x) \;=\; \int \mathrm{d}a\; P(a,x),
\end{equation}
for the conditional. The second equation --- integrating out what one does not care about ---
is called \emph{marginalisation}, and it is the operation that will turn out to be expensive
and that all the machinery in Chapter~\ref{ch:bp} exists to make cheap.

Bayes' rule is the same statement rearranged:
\begin{equation}
\label{eq:bayes}
P(a \mid x) \;=\; \frac{P(x\mid a)\,P(a)}{P(x)} \;\propto\; \underbrace{P(x\mid a)}_{\text{likelihood}}\;\underbrace{P(a)}_{\text{prior}} .
\end{equation}
The proportionality is in $a$: the denominator $P(x)$ does not depend on $a$ and serves only
to normalise. Almost every calculation in this document works with the unnormalised product
and restores the constant at the end, because the constant is usually the expensive part.

\begin{intuition}
The prior says what clean data look like before seeing anything. The likelihood says how
plausible the observation is for each candidate clean value. Their product, renormalised, is
what one should believe after seeing the observation. In our setting the prior is the Markov
chain and the likelihood is the noise model.
\end{intuition}

\section{Expectation and conditional expectation}

The expectation of $g(a)$ is $\E[g(a)] = \int \mathrm{d}a\, P(a) g(a)$. The
\emph{conditional} expectation given $x$ is the same integral against the conditional density,
\begin{equation}
\E[g(a)\mid x] \;=\; \int \mathrm{d}a\; P(a\mid x)\, g(a).
\end{equation}
Following the notation of the diffusion literature we are matching, we write the conditional
mean of the clean signal with angle brackets,
\begin{equation}
\pmean{a} \;:=\; \E[a \mid x, t] \;=\; \int \mathrm{d}a\; P(a\mid x,t)\, a ,
\end{equation}
where the subscript records that the conditioning is on the observation $x$ at noise level
$t$. Componentwise, $\pmeansite{i}$ is the posterior mean at site $i$.

A fact used repeatedly: conditional expectation is a \emph{function of $x$}, not a number. It
is the function that will be estimated, whether by exact computation or by a neural network.

\section{The Bayes-optimal estimator under squared loss}
\label{sec:bayesopt}

This is the single most-used result in the document, so it is proved rather than quoted.

\begin{proposition}[Orthogonality]
\label{prop:orthogonality}
Let $a$ have finite second moment. For any $u$ that is a function of $x$ alone,
\begin{equation}
\label{eq:orth}
\E\!\left[\|a-u\|^2 \,\middle|\, x\right]
\;=\;
\E\!\left[\|a-\pmean{a}\|^2 \,\middle|\, x\right]
\;+\;
\|u-\pmean{a}\|^2 .
\end{equation}
Consequently $\pmean{a}$ minimises the conditional mean squared error, and hence also the
unconditional one.
\end{proposition}

\begin{proof}
Insert and subtract $\pmean{a}$:
\[
\|a-u\|^2 = \|(a - \pmean{a}) + (\pmean{a} - u)\|^2
= \|a-\pmean{a}\|^2 + 2\langle a - \pmean{a},\, \pmean{a}-u\rangle + \|\pmean{a}-u\|^2 .
\]
Take $\E[\cdot\mid x]$. In the cross term, $\pmean{a}-u$ is a function of $x$ only, so it is
constant under the conditional expectation and can be pulled out:
\[
\E\!\left[\langle a - \pmean{a},\, \pmean{a}-u\rangle \,\middle|\, x\right]
= \big\langle \E[a\mid x] - \pmean{a},\; \pmean{a}-u \big\rangle = 0 ,
\]
since $\E[a\mid x] = \pmean{a}$ by definition. The remaining two terms give \eqref{eq:orth}.
The first is independent of $u$ and the second is non-negative and vanishes only at
$u = \pmean{a}$. Taking a further expectation over $x$ preserves the minimiser.
\end{proof}

Two consequences deserve emphasis.

\paragraph{There is an irreducible floor.} The first term of \eqref{eq:orth} does not depend
on the estimator. It is the posterior variance, and no denoiser --- exact, learned, or
otherwise --- can go below it. A reported denoising loss is therefore uninterpretable without
knowing that floor.

\begin{pitfall}
Early in this project a test asserted that a network should achieve a $40\%$ reduction in
denoising loss. That target was arithmetically impossible: with the floor at $7.4$ and the
loss of the trivial predictor $f\equiv0$ at $12.06$, the largest achievable reduction was
$39\%$. The test could never pass however well the network trained. The fix was to report the
\emph{fraction of the achievable reduction} rather than the raw one.
\begin{codebox}
\coderef{src/sample\_metrics.py}{bayes\_risk} computes the floor;
\coderef{src/sample\_metrics.py}{pointwise\_ladder} returns
\texttt{fraction\_achievable\_captured} alongside the raw loss for exactly this reason.
\end{codebox}
\end{pitfall}

\paragraph{Optimality is loss-specific.} Proposition~\ref{prop:orthogonality} is about squared
error. Under absolute error the optimal estimator is the posterior median; under $0$--$1$ loss
it is the mode. Diffusion training uses squared error, which is why the posterior \emph{mean}
is the relevant functional --- not because it is intrinsically the right summary of a
posterior.

\section{Linear estimators and LMMSE}
\label{sec:lmmse-intro}

Sometimes one restricts to estimators of the form $\widehat{a} = Bx$ for a matrix $B$. The
best such estimator under squared loss is the \emph{linear minimum mean squared error}
estimator, and it depends on the joint distribution only through second moments.

\begin{proposition}
\label{prop:lmmse-general}
Let $a, x$ have zero mean and finite second moments, with $\Cov(x)$ invertible. Then
\begin{equation}
\label{eq:lmmse-general}
\lmmse \;=\; \Cov(a,x)\,\Cov(x)^{-1} x
\end{equation}
minimises $\E\|a - Bx\|^2$ over matrices $B$.
\end{proposition}

\begin{proof}
Expand $\E\|a-Bx\|^2 = \operatorname{tr}\!\big(\Cov(a) - 2B\Cov(x,a) + B\Cov(x)B^{\!\top}\big)$
and differentiate with respect to $B$:
$-2\Cov(a,x) + 2B\Cov(x) = 0$, so $B = \Cov(a,x)\Cov(x)^{-1}$. The objective is convex in $B$
because $\Cov(x)\succ0$, so the stationary point is the minimum.
\end{proof}

\begin{intuition}
LMMSE is what one gets from knowing only the covariance structure. If the joint distribution
happens to be Gaussian, the conditional mean is already linear and LMMSE coincides with the
Bayes-optimal estimator. Otherwise LMMSE is strictly worse, and the gap is exactly the value
of the higher-order structure. Chapter~\ref{ch:gaussian} shows that a natural and widely used
approximation to exact inference silently computes LMMSE, which is why the distinction
matters so much here.
\end{intuition}

\section{Two ways a distribution can be ``the data''}
\label{sec:empirical}

This distinction causes a common error, so it is fixed now.

The \emph{population} distribution $\Pz$ is the true source of the data. One never has it ---
one has $\nseq$ samples $a^1,\dots,a^{\nseq}$ drawn from it. Those samples define the
\emph{empirical} distribution
\begin{equation}
\label{eq:empirical}
\widehat{P}_0 \;=\; \frac{1}{\nseq}\sum_{\mu=1}^{\nseq} \delta_{a^\mu},
\end{equation}
a sum of point masses. It is a perfectly good probability distribution, and it is not the
population one.

Correspondingly there are two versions of any quantity defined from a distribution: the
population posterior mean, computed under $\Pz$, and the empirical posterior mean, computed
under $\widehat{P}_0$. A training procedure that minimises an average over the observed
samples is minimising the \emph{empirical} objective, and its unrestricted optimum is the
\emph{empirical} posterior mean.

\begin{pitfall}
It is tempting --- and wrong --- to say that the population objective is minimised by the
empirical score. The population objective is minimised by the population posterior mean, by
Proposition~\ref{prop:orthogonality}; it is the \emph{empirical} objective whose optimum is
the empirical posterior mean. The distinction is the whole content of the memorisation
discussion in \S\ref{sec:memorisation}, and conflating the two makes that discussion
incoherent. An earlier draft of the accompanying note contained exactly this error.
\end{pitfall}

\section*{Sources}
\addcontentsline{toc}{section}{Sources}
Conditional expectation as the minimiser of squared error, and the orthogonality principle
behind it, are standard: see \citet{lehmann1998} for the estimation-theoretic treatment and
\citet{kay1993} for the signal-processing one, where the linear case is developed as the LMMSE
estimator used in Chapter~\ref{ch:gaussian}.
